% GENERATED FILE. Edit canonical lessons, then run npm run generate:books.
% canonical-source-hash: fnv1a64:8661e99a37f327d9
% canonical-lessons: TA-C325-olaiccuvati, TA-C325-kalvettu, TA-C325-pulavar, TA-C325-kaviyam, TA-C325-iyarru, TA-R325-review-318-325

\chapter{Old Writings}
\label{ch:ta-old-writings}
\hlchaptermodality{325}

\begin{chapteropening}
\textbf{By the end of this chapter:} \emph{I can talk about old manuscripts, inscriptions and poets.}

Complete the last lesson of chapter 325: I can talk about old manuscripts, inscriptions and poets.
\end{chapteropening}

\section[\texorpdfstring{\ta{ஓலைச்சுவடி} (ōlaiccuvaṭi)}{ஓலைச்சுவடி (ōlaiccuvaṭi)}]{\ta{ஓலைச்சுவடி} (ōlaiccuvaṭi) — a palm-leaf manuscript}

\label{lesson:TA-C325-olaiccuvati}



\begin{quote}
Before the new one: say the Tamil for to cremate, then the Tamil for a memorial meeting.
\end{quote}

\subsection*{You'll want to know: \ta{ஓலைச்சுவடி}}
\textbf{\ta{ஓலைச்சுவடி}} — \emph{ōlaiccuvaṭi} — \textquotedblleft{}a palm-leaf manuscript\textquotedblright{}.

\begin{quote}
\ta{ஈரம்} \ta{பட்டால்} \ta{ஓலைச்சுவடி} \ta{கெட்டுவிடும்}.
\end{quote}

\begin{quote}
\emph{īram paṭṭāl ōlaiccuvaṭi keṭṭuviṭum.}
\end{quote}

\begin{quote}
If damp gets to it, a palm-leaf manuscript is ruined.
\end{quote}

Plural \ta{ஓலைச்சுவடிகள்}; \ta{ஓலை}, palm leaf + \ta{சுவடி}, a bundle of leaves.

\begin{cousinweb}
\textbf{\ta{ஓலைச்சுவடி}} (\emph{ōlaiccuvaṭi}) is a Tamil compound built from \ta{ஓலை}, \textquotedblleft{}palm leaf, writing on a palm leaf\textquotedblright{} (DEDR 1070), and \ta{சுவடி}, \textquotedblleft{}a bound bundle of leaves\textquotedblright{}, whose origin is not settled.
\end{cousinweb}

\begin{grammarlens}[title={What you've built}]
One more word for telling what happened.
\end{grammarlens}

\subsection*{Your turn}
\begin{itemize}
\raggedright
  \item \emph{Say it:} \emph{ōlaiccuvaṭi}
  \item \emph{Say it:} \emph{ōlaiccuvaṭi}, once more
  \item \emph{Say it:} \emph{īram paṭṭāl ōlaiccuvaṭi keṭṭuviṭum.} — If damp gets to it, a palm-leaf manuscript is ruined.
  \item \emph{Recall:} say the Tamil for to cremate, then the Tamil for a memorial meeting, then say \emph{ōlaiccuvaṭi} again
\end{itemize}

\begin{tcolorbox}[breakable,colback=teal!4,colframe=teal!35!black,title={Before you move on}]
What does \ta{ஓலைச்சுவடி} mean? (\textquotedblleft{}A palm-leaf manuscript\textquotedblright{}.) Say it once more.
\end{tcolorbox}
\section[\texorpdfstring{\ta{கல்வெட்டு} (kalveṭṭu)}{கல்வெட்டு (kalveṭṭu)}]{\ta{கல்வெட்டு} (kalveṭṭu) — an inscription}

\label{lesson:TA-C325-kalvettu}



\begin{quote}
Before the new one: say the Tamil for a memorial meeting, then the Tamil for a palm-leaf manuscript.
\end{quote}

\subsection*{You'll want to know: \ta{கல்வெட்டு}}
\textbf{\ta{கல்வெட்டு}} — \emph{kalveṭṭu} — \textquotedblleft{}an inscription\textquotedblright{}.

\begin{quote}
\ta{கோயில்} \ta{சுவரில்} \ta{ஒரு} \ta{பழைய} \ta{கல்வெட்டு} \ta{இருக்கிறது}.
\end{quote}

\begin{quote}
\emph{kōyil cuvaril oru paḻaiya kalveṭṭu irukkiṟatu.}
\end{quote}

\begin{quote}
There is an old inscription on the temple wall.
\end{quote}

\ta{கல்}, stone + \ta{வெட்டு}, to cut: writing cut in stone.

\begin{cousinweb}
\textbf{\ta{கல்வெட்டு}} (\emph{kalveṭṭu}) is a Tamil compound built from \ta{கல்}, \textquotedblleft{}stone\textquotedblright{} (DEDR 1298), and \ta{வெட்டு}, \textquotedblleft{}to cut, engrave\textquotedblright{} (DEDR 5478).
\end{cousinweb}

\begin{grammarlens}[title={What you've built}]
One more word for telling what happened.
\end{grammarlens}

\subsection*{Your turn}
\begin{itemize}
\raggedright
  \item \emph{Say it:} \emph{kalveṭṭu}
  \item \emph{Say it:} \emph{kalveṭṭu}, once more
  \item \emph{Say it:} \emph{kōyil cuvaril oru paḻaiya kalveṭṭu irukkiṟatu.} — There is an old inscription on the temple wall.
  \item \emph{Recall:} say the Tamil for a memorial meeting, then the Tamil for a palm-leaf manuscript, then say \emph{kalveṭṭu} again
\end{itemize}

\begin{tcolorbox}[breakable,colback=teal!4,colframe=teal!35!black,title={Before you move on}]
What does \ta{கல்வெட்டு} mean? (\textquotedblleft{}An inscription\textquotedblright{}.) Say it once more.
\end{tcolorbox}
\section[\texorpdfstring{\ta{புலவர்} (pulavar)}{புலவர் (pulavar)}]{\ta{புலவர்} (pulavar) — a classical poet; a scholar}

\label{lesson:TA-C325-pulavar}



\begin{quote}
Before the new one: say the Tamil for a palm-leaf manuscript, then the Tamil for an inscription.
\end{quote}

\subsection*{You'll want to know: \ta{புலவர்}}
\textbf{\ta{புலவர்}} — \emph{pulavar} — \textquotedblleft{}a classical poet; a scholar\textquotedblright{}.

\begin{quote}
\ta{அந்தப்} \ta{புலவர்} \ta{பல} \ta{பாடல்களை} \ta{எழுதினார்}.
\end{quote}

\begin{quote}
\emph{antap pulavar pala pāṭalkaḷai eḻutiṉār.}
\end{quote}

\begin{quote}
That poet wrote many poems.
\end{quote}

Said of the classical poets and of Tamil scholars.

\begin{cousinweb}
\textbf{\ta{புலவர்}} (\emph{pulavar}), \textquotedblleft{}a learned poet, a scholar\textquotedblright{}, is a Tamil word whose origin is not settled in the sources at hand.
\end{cousinweb}

\begin{grammarlens}[title={What you've built}]
One more word for telling what happened.
\end{grammarlens}

\subsection*{Your turn}
\begin{itemize}
\raggedright
  \item \emph{Say it:} \emph{pulavar}
  \item \emph{Say it:} \emph{pulavar}, once more
  \item \emph{Say it:} \emph{antap pulavar pala pāṭalkaḷai eḻutiṉār.} — That poet wrote many poems.
  \item \emph{Recall:} say the Tamil for a palm-leaf manuscript, then the Tamil for an inscription, then say \emph{pulavar} again
\end{itemize}

\begin{tcolorbox}[breakable,colback=teal!4,colframe=teal!35!black,title={Before you move on}]
What does \ta{புலவர்} mean? (\textquotedblleft{}A classical poet; a scholar\textquotedblright{}.) Say it once more.
\end{tcolorbox}
\section[\texorpdfstring{\ta{காவியம்} (kāviyam)}{காவியம் (kāviyam)}]{\ta{காவியம்} (kāviyam) — an epic}

\label{lesson:TA-C325-kaviyam}



\begin{quote}
Before the new one: say the Tamil for an inscription, then the Tamil for a classical poet.
\end{quote}

\subsection*{You'll want to know: \ta{காவியம்}}
\textbf{\ta{காவியம்}} — \emph{kāviyam} — \textquotedblleft{}an epic\textquotedblright{}.

\begin{quote}
\ta{இந்தக்} \ta{காவியம்} \ta{ஒரு} \ta{வீரனின்} \ta{கதை}.
\end{quote}

\begin{quote}
\emph{intak kāviyam oru vīraṉiṉ katai.}
\end{quote}

\begin{quote}
This epic is the story of a hero.
\end{quote}

Plural \ta{காவியங்கள்}.

\begin{cousinweb}
\textbf{\ta{காவியம்}} (\emph{kāviyam}) is borrowed from Sanskrit \emph{kāvya}, \textquotedblleft{}poetical, a poem\textquotedblright{} (Monier-Williams).
\end{cousinweb}

\begin{grammarlens}[title={What you've built}]
One more word for telling what happened.
\end{grammarlens}

\subsection*{Your turn}
\begin{itemize}
\raggedright
  \item \emph{Say it:} \emph{kāviyam}
  \item \emph{Say it:} \emph{kāviyam}, once more
  \item \emph{Say it:} \emph{intak kāviyam oru vīraṉiṉ katai.} — This epic is the story of a hero.
  \item \emph{Recall:} say the Tamil for an inscription, then the Tamil for a classical poet, then say \emph{kāviyam} again
\end{itemize}

\begin{tcolorbox}[breakable,colback=teal!4,colframe=teal!35!black,title={Before you move on}]
What does \ta{காவியம்} mean? (\textquotedblleft{}An epic\textquotedblright{}.) Say it once more.
\end{tcolorbox}
\section[\texorpdfstring{\ta{இயற்று} (iyaṟṟu)}{இயற்று (iyaṟṟu)}]{\ta{இயற்று} (iyaṟṟu) — to compose}

\label{lesson:TA-C325-iyarru}



\begin{quote}
Before the new one: say the Tamil for a classical poet, then the Tamil for an epic.
\end{quote}

\subsection*{You'll want to know: \ta{இயற்று}}
\textbf{\ta{இயற்று}} — \emph{iyaṟṟu} — \textquotedblleft{}to compose\textquotedblright{}.

\begin{quote}
\ta{அவர்} \ta{இந்த} \ta{நூலை} \ta{இயற்றினார்}.
\end{quote}

\begin{quote}
\emph{avar inta nūlai iyaṟṟiṉār.}
\end{quote}

\begin{quote}
He composed this work.
\end{quote}

Past (I …): \ta{இயற்றினேன்}, \emph{iyaṟṟiṉēṉ}.

Of composing a poem or a book, in written Tamil.

\begin{cousinweb}
\textbf{\ta{இயற்று}} (\emph{iyaṟṟu}) is native Dravidian: DEDR 471 gives it beside \ta{இயல்}, \textquotedblleft{}to be possible\textquotedblright{}, as \textquotedblleft{}to do, effect, cause to act\textquotedblright{}.
\end{cousinweb}

\begin{grammarlens}[title={What you've built}]
One more word for telling what happened.
\end{grammarlens}

\subsection*{Your turn}
\begin{itemize}
\raggedright
  \item \emph{Say it:} \emph{iyaṟṟu}
  \item \emph{Say it:} \emph{iyaṟṟu}, once more
  \item \emph{Say it:} \emph{avar inta nūlai iyaṟṟiṉār.} — He composed this work.
  \item \emph{Recall:} say the Tamil for a classical poet, then the Tamil for an epic, then say \emph{iyaṟṟu} again
\end{itemize}

\begin{tcolorbox}[breakable,colback=teal!4,colframe=teal!35!black,title={Before you move on}]
What does \ta{இயற்று} mean? (\textquotedblleft{}To compose\textquotedblright{}.) Say it once more.
\end{tcolorbox}
\section[(dialogue)]{Looking back: from An Opening Ceremony to Old Writings}

\label{lesson:TA-R325-review-318-325}



\begin{quote}
Say the Tamil for an opening ceremony and a standing lamp.
\end{quote}

\subsection*{Your turn}
Say these aloud:

\begin{itemize}
\raggedright
  \item an opening ceremony, a standing lamp, a chief guest, to preside, to give a speech
  \item an anniversary, a golden jubilee, an annual day, a festive look, an invitation card
  \item to hoist a flag, a parade, a festoon, a banner, to proclaim
  \item a ritual, a convention, to name a child, a baby shower, a housewarming
  \item to wed, a wedding pendant, an auspicious hour, a wedding dais, a pandal
  \item a funeral, a procession, a wreath, to bury, a grave
  \item a tribute, last respects, a memorial day, to cremate, a memorial meeting
  \item a palm-leaf manuscript, an inscription, a classical poet, an epic, to compose
\end{itemize}

\begin{tcolorbox}[breakable,colback=teal!4,colframe=teal!35!black,title={Before you move on}]
An opening ceremony? (\textbf{\ta{திறப்புவிழா}}, \emph{tiṟappuviḻā}.) A standing lamp? (\textbf{\ta{குத்துவிளக்கு}}, \emph{kuttuviḷakku}.) A chief guest? (\textbf{\ta{சிறப்பு} \ta{விருந்தினர்}}, \emph{ciṟappu viruntiṉar}.) To preside? (\textbf{\ta{தலைமைதாங்கு}}, \emph{talaimaitāṅku}.) To give a speech? (\textbf{\ta{உரையாற்று}}, \emph{uraiyāṟṟu}.) An anniversary? (\textbf{\ta{ஆண்டு} \ta{நிறைவு}}, \emph{āṇṭu niṟaivu}.) A golden jubilee? (\textbf{\ta{பொன்விழா}}, \emph{poṉviḻā}.) An annual day? (\textbf{\ta{ஆண்டுவிழா}}, \emph{āṇṭuviḻā}.) A festive look? (\textbf{\ta{விழாக்கோலம்}}, \emph{viḻākkōlam}.) An invitation card? (\textbf{\ta{அழைப்பிதழ்}}, \emph{aḻaippitaḻ}.) To hoist a flag? (\textbf{\ta{கொடியேற்று}}, \emph{koṭiyēṟṟu}.) A parade? (\textbf{\ta{அணிவகுப்பு}}, \emph{aṇivakuppu}.) A festoon? (\textbf{\ta{தோரணம்}}, \emph{tōraṇam}.) A banner? (\textbf{\ta{பதாகை}}, \emph{patākai}.) To proclaim? (\textbf{\ta{பறைசாற்று}}, \emph{paṟaicāṟṟu}.) A ritual? (\textbf{\ta{சடங்கு}}, \emph{caṭaṅku}.) A convention? (\textbf{\ta{மரபு}}, \emph{marapu}.) To name a child? (\textbf{\ta{பெயர்சூட்டு}}, \emph{peyarcūṭṭu}.) A baby shower? (\textbf{\ta{வளைகாப்பு}}, \emph{vaḷaikāppu}.) A housewarming? (\textbf{\ta{புதுமனை} \ta{புகுவிழா}}, \emph{putumaṉai pukuviḻā}.) To wed? (\textbf{\ta{மணமுடி}}, \emph{maṇamuṭi}.) A wedding pendant? (\textbf{\ta{தாலி}}, \emph{tāli}.) An auspicious hour? (\textbf{\ta{முகூர்த்தம்}}, \emph{mukūrttam}.) A wedding dais? (\textbf{\ta{மணமேடை}}, \emph{maṇamēṭai}.) A pandal? (\textbf{\ta{பந்தல்}}, \emph{pantal}.) A funeral? (\textbf{\ta{இறுதிச்சடங்கு}}, \emph{iṟuticcaṭaṅku}.) A procession? (\textbf{\ta{ஊர்வலம்}}, \emph{ūrvalam}.) A wreath? (\textbf{\ta{மலர்வளையம்}}, \emph{malarvaḷaiyam}.) To bury? (\textbf{\ta{புதை}}, \emph{putai}.) A grave? (\textbf{\ta{கல்லறை}}, \emph{kallaṟai}.) A tribute? (\textbf{\ta{அஞ்சலி}}, \emph{añcali}.) Last respects? (\textbf{\ta{இறுதிமரியாதை}}, \emph{iṟutimariyātai}.) A memorial day? (\textbf{\ta{நினைவுநாள்}}, \emph{niṉaivunāḷ}.) To cremate? (\textbf{\ta{எரியூட்டு}}, \emph{eriyūṭṭu}.) A memorial meeting? (\textbf{\ta{நினைவேந்தல்}}, \emph{niṉaivēntal}.) A palm-leaf manuscript? (\textbf{\ta{ஓலைச்சுவடி}}, \emph{ōlaiccuvaṭi}.) An inscription? (\textbf{\ta{கல்வெட்டு}}, \emph{kalveṭṭu}.) A classical poet? (\textbf{\ta{புலவர்}}, \emph{pulavar}.) An epic? (\textbf{\ta{காவியம்}}, \emph{kāviyam}.) To compose? (\textbf{\ta{இயற்று}}, \emph{iyaṟṟu}.)
\end{tcolorbox}
